\chapter{Sizes and Their Precedences}

APFS records the size of a file in several places, and the places do
not always agree. A parser that picks the wrong source on the wrong
file silently produces wrong rows and wrong directory totals. This
chapter fixes two rules --- one for logical size, one for allocated
size --- and shows the cases for which each is the correct rule.

The two metrics are independent. They have separate sources, separate
oracles, and separate fail-closed boundaries. The chapter develops
them in that order: logical first (the primary anchor; covers every
case), allocated second (covers five of seven case classes; sparse
and decmpfs fail closed).

\section{Logical Size}

The primary size metric is \emph{logical size}: the user-visible byte
length, the number POSIX \texttt{st\_size} reports, the number the
shell shows in \texttt{ls -l}. It is not the bytes on disk; it is
not the bytes the file would occupy after decompression on a
different filesystem; it is not the bytes uniquely owned by this
inode. Each of those is a separate metric, and at most one is
treated here (allocated size, in the next section).

The choice of logical size as the primary anchor is deliberate. It
has a stable oracle (\texttt{st\_size}), applies uniformly to
ordinary, sparse, cloned, hard-linked, symlink, and compressed
files, and is the size every user-facing tool already agrees on.

\subsection{The Precedence Rule}

\textbf{Observation.} Five steps, applied per inode, produce
\texttt{st\_size} for every case the validated corpus covers:

\begin{enumerate}
  \item If the inode's \texttt{internal\_flags} carry
    \texttt{INODE\_HAS\_UNCOMPRESSED\_SIZE}: pick
    \texttt{inode.uncompressed\_size}.
  \item Else if the inode carries a \texttt{j\_dstream\_t} (via
    \texttt{INO\_EXT\_TYPE\_DSTREAM} or a paired
    \texttt{j\_dstream\_id\_val\_t} record): pick
    \texttt{j\_dstream\_t.size}.
  \item Else if a \texttt{com.apple.decmpfs} xattr is present and
    step 1 did not fire: pick the decmpfs header's
    \texttt{uncompressed\_size} field.
  \item For symlinks (\texttt{S\_IFLNK}): pick the byte length of the
    \texttt{com.apple.fs.symlink} xattr payload.
  \item Otherwise: zero.
\end{enumerate}

The order matters. Steps 1 and 3 both reach into compression
metadata, but they reach for different fields and have different
trust. The next subsection shows why.

\begin{figure}[h]
\centering
\begin{tikzpicture}[
  every node/.style={font=\footnotesize},
  decision/.style={diamond, draw, aspect=2.4, inner sep=1pt,
                   align=center, fill=blue!8},
  result/.style={draw, rounded corners=2pt, fill=green!18,
                 minimum width=3.6cm, minimum height=0.8cm, align=center},
  fallback/.style={draw, rounded corners=2pt, fill=gray!15,
                   minimum width=3.0cm, minimum height=0.8cm, align=center},
  flow/.style={-Stealth},
  yeslab/.style={font=\scriptsize, midway, above, sloped=false,
                 fill=white, inner sep=1pt},
  nolab/.style={font=\scriptsize, midway, right=2pt,
                fill=white, inner sep=1pt},
]
  \node[decision] (d1) {INODE\_HAS\_\\UNCOMPRESSED\_SIZE?};
  \node[result, right=2.0cm of d1] (r1) {\texttt{inode.uncompressed\_size}};

  \node[decision, below=1.2cm of d1] (d2) {\texttt{j\_dstream\_t}\\present?};
  \node[result, right=2.0cm of d2] (r2) {\texttt{j\_dstream\_t.size}};

  \node[decision, below=1.2cm of d2] (d3) {decmpfs xattr\\present?};
  \node[result, right=2.0cm of d3] (r3) {decmpfs header\\\texttt{uncompressed\_size}};

  \node[decision, below=1.2cm of d3] (d4) {symlink\\(S\_IFLNK)?};
  \node[result, right=2.0cm of d4] (r4) {symlink target\\byte length};

  \node[fallback, below=1.2cm of d4] (r5) {zero};

  \draw[flow] (d1.east) -- (r1.west) node[yeslab] {yes};
  \draw[flow] (d2.east) -- (r2.west) node[yeslab] {yes};
  \draw[flow] (d3.east) -- (r3.west) node[yeslab] {yes};
  \draw[flow] (d4.east) -- (r4.west) node[yeslab] {yes};

  \draw[flow] (d1.south) -- (d2.north) node[nolab] {no};
  \draw[flow] (d2.south) -- (d3.north) node[nolab] {no};
  \draw[flow] (d3.south) -- (d4.north) node[nolab] {no};
  \draw[flow] (d4.south) -- (r5.north) node[nolab] {no};
\end{tikzpicture}
\caption{Per-inode logical-size precedence. The order is non-negotiable;
  the compressed case is the load-bearing reason step~1 must beat
  step~3.}
\end{figure}

In pseudocode the same rule reads as one short function:

\begin{minted}{python}
def logical_size(inode, xattrs):
    if inode.internal_flags & INODE_HAS_UNCOMPRESSED_SIZE:
        return inode.uncompressed_size
    if inode.dstream is not None:
        return inode.dstream.size
    decmpfs = xattrs.get("com.apple.decmpfs")
    if decmpfs is not None:
        return decmpfs.header.uncompressed_size
    if S_ISLNK(inode.mode):
        return len(xattrs["com.apple.fs.symlink"].payload)
    return 0
\end{minted}

\subsection{The Load-Bearing Case: Compression}

\textbf{Observation.} On the EX-19 fixture the compressed file
(\texttt{compressed.txt}, created via
\texttt{ditto --hfsCompression}) presents:

\begin{itemize}
  \item \texttt{INODE\_HAS\_UNCOMPRESSED\_SIZE} set,
    \texttt{inode.uncompressed\_size = 53248}.
  \item No \texttt{j\_dstream\_t}. APFS stores small compressed data
    inline in the decmpfs xattr and does not allocate a separate
    dstream for it.
  \item \texttt{com.apple.decmpfs} xattr present, with header
    \texttt{uncompressed\_size = 437} and
    \texttt{compression\_type = 0}.
  \item Mounted POSIX \texttt{st\_size = 53248}.
\end{itemize}

The decmpfs header's \texttt{uncompressed\_size} is a placeholder for
this layout. Step~1 must beat step~3, or the parser emits 437 instead
of 53248. The fixture is small but the principle is general: when
APFS sets the inode flag, the inode flag is the source of truth;
the decmpfs header is metadata the running system maintains for its
own decompression path and is not reliably the file's size.

\subsection{The Hint That Is Not A Size}

\textbf{Observation.}
\texttt{INO\_EXT\_TYPE\_SPARSE\_BYTES} is an
\emph{allocation hint}, not a logical size. It records the bytes the
sparse-allocation layer has not committed to disk.

On the proof fixture's sparse file:

\begin{itemize}
  \item \texttt{j\_dstream\_t.size = 1052897} (which equals
    \texttt{st\_size}).
  \item \texttt{INO\_EXT\_TYPE\_SPARSE\_BYTES = 1032192}.
  \item \texttt{st\_size = 1052897}.
\end{itemize}

The sparse bytes number describes the difference between the file's
logical extent and its allocated bytes. It is useful for the
allocated-size metric (next section) but is not the answer to "how
large is this file"; that number is the dstream size, which the
precedence rule correctly picks at step~2.

A reader that treats \texttt{SPARSE\_BYTES} as anything else is
adding sparse holes to the logical size of every sparse file --- a
sizeable error on real-world data.

\subsection{Per-Inode Results}

\textbf{Observation.} On the EX-19 same-run fixture covering ordinary,
sparse, clone, hard-linked, symlink, and compressed cases, the
precedence rule reproduces \texttt{st\_size} for every unique inode:

\begin{center}
\begin{tabular}{lrlrr}
\toprule
path & inode & step picked & value & st\_size \\
\midrule
ordinary.txt   & 16 & \texttt{j\_dstream\_size}         & 29      & 29 \\
sparse.bin     & 19 & \texttt{j\_dstream\_size}         & 1052897 & 1052897 \\
clone.txt      & 20 & \texttt{j\_dstream\_size}         & 29      & 29 \\
link.txt       & 23 & \texttt{symlink\_target\_len}     & 12      & 12 \\
compressed.txt & 24 & \texttt{inode\_uncompressed\_size} & 53248   & 53248 \\
\bottomrule
\end{tabular}
\end{center}

Zero mismatches. The compressed case picks step~1; everything else
picks step~2 except the symlink (step~4). The hard link does not
appear as a separate row because the rule is per inode and
\texttt{hard.txt} shares inode~16 with \texttt{ordinary.txt}.

\section{Allocated Size}

Per inode, the number of bytes APFS has allocated on disk for the
file's content stream.

\textbf{Observation.} For ordinary, cloned, hard-linked, and symlink
cases the parser emits \texttt{j\_dstream\_t.alloced\_size} (or zero
for symlinks) and the result matches macOS's \texttt{st\_blocks *
512} oracle. Two further cases used to fail closed --- sparse files
and decmpfs-compressed files --- and were closed by EX-26 against
the same \texttt{st\_blocks * 512} oracle. For sparse, the
allocated bytes are \texttt{alloced\_size - sparse\_bytes}. For
decmpfs, the allocated bytes are the sum of any stream-backed
xattrs' \texttt{alloced\_size} values (the primary dstream is
usually absent, and the compressed bytes live in either
\texttt{com.apple.decmpfs} or \texttt{com.apple.ResourceFork}; an
embedded xattr contributes zero because the bytes live inline in
the xattr payload).

\subsection{The Precedence Rule}

\textbf{Observation.} Six branches, applied per inode:

\begin{enumerate}
  \item Regular file with a \texttt{j\_dstream\_t} xfield and no
    \texttt{INO\_EXT\_TYPE\_SPARSE\_BYTES} xfield and no
    \texttt{com.apple.decmpfs} xattr: emit
    \texttt{Some(j\_dstream\_t.alloced\_size)}.
  \item Regular file with a \texttt{j\_dstream\_t} xfield and a
    \texttt{SPARSE\_BYTES} xfield (no decmpfs): emit
    \texttt{Some(alloced\_size - sparse\_bytes)} (EX-26 Hypothesis~A,
    see §\ref{ssec:sparse-bytes-lift}).
  \item Regular file with a \texttt{com.apple.decmpfs} xattr: emit
    \texttt{Some(primary + decmpfs.stream\_dstream.alloced +
    ResourceFork.stream\_dstream.alloced)} (each term defaults to
    zero when absent or embedded; EX-26 Hypothesis~F, see
    §\ref{ssec:decmpfs-lift}).
  \item Symlinks: emit \texttt{Some(0)}.
  \item Directories: emit \texttt{Some(0)}.
  \item Anything else: emit \texttt{None}.
\end{enumerate}

\texttt{None} is not a synonym for zero. A row with
\texttt{allocated\_size = None} is recorded in the run's
\texttt{not\_claimed} register, the metric toggle in the
visualisation greys out for that row, and any directory aggregate
that contains a \texttt{None}-allocated inode is itself emitted as
\texttt{None} --- a partial aggregate is forbidden because a
consumer might otherwise treat it as authoritative.

\subsection{Why \texttt{alloced\_size} Is The Only Candidate}

Six independent open-source readers were surveyed
(\texttt{linux-apfs-rw}, \texttt{apfsprogs}/\texttt{apfsck},
the Sleuth Kit, \texttt{libfsapfs}, \texttt{dissect.apfs},
\texttt{apfs-fuse}, \texttt{go-apfs}). All six decode
\texttt{j\_dstream\_t.alloced\_size}; only TSK and
\texttt{linux-apfs-rw} surface it to callers, and
\texttt{linux-apfs-rw} publishes it directly as
\texttt{inode->i\_blocks << 9}. No reader surfaces the alternative
candidate --- the sum of \texttt{j\_file\_extent\_val\_t.len} over
the inode's extents --- as a public allocated-bytes number. The
extent-reference tree
(\texttt{OBJECT\_TYPE\_BLOCKREFTREE}) is opened by some readers but
is the per-volume reference-counting authority for snapshots and
clones, not a per-file allocated-bytes table.

A consistency conflict surfaces between two halves of the same
project: \texttt{linux-apfs-rw} (the kernel module) writes
\texttt{alloced\_size = round\_up(ds\_size, blocksize)};
\texttt{apfsck} (the same project's filesystem checker) enforces
\texttt{alloced\_size == sum(j\_file\_extent\_val\_t.len)}. The two
agree only when every extent ends exactly at \texttt{ds\_size}
and the file has no truncate residue. The kernel can therefore
write volumes that its companion checker flags as corrupt. The
right oracle is the external public number
(\texttt{st\_blocks * 512}), not internal field consistency.

\subsection{The Sparse-File Lift (EX-26 Hypothesis A)}
\label{ssec:sparse-bytes-lift}

EX-22 left this branch fail-closed because the fixture carried a
single sparse shape and the observed identity
\texttt{alloced\_size - sparse\_bytes = st\_blocks * 512} was a
hypothesis worth a corpus-level probe before promoting to rule.
EX-26 ran that probe on four sparse shapes and the identity held
exactly across all of them.

\textbf{Observation.} On the EX-22 fixture the sparse file carries:

\begin{itemize}
  \item \texttt{j\_dstream\_t.size = 1052897} (equals \texttt{st\_size}).
  \item \texttt{j\_dstream\_t.alloced\_size = 1056768}
    ($= \text{round\_up}(\text{size}, 4096)$).
  \item \texttt{INO\_EXT\_TYPE\_SPARSE\_BYTES = 1032192}
    (the unallocated bytes inside the logical extent).
  \item \texttt{st\_blocks = 48}, so
    \texttt{st\_blocks * 512 = 24576}.
\end{itemize}

\texttt{alloced\_size} overstates the actual on-disk allocation by
exactly $1{,}056{,}768 - 24{,}576 = 1{,}032{,}192$ bytes --- the
sparse-bytes hint. APFS on macOS follows the
\texttt{linux-apfs-rw} kernel rule rather than the
\texttt{apfsck} sum-of-extents rule. The emission rule must
therefore split the dstream case by presence of the
\texttt{SPARSE\_BYTES} xfield.

\textbf{EX-26 corpus result.} The identity
\texttt{alloced\_size - sparse\_bytes = st\_blocks * 512} holds
exactly across four sparse shapes:

\begin{center}
\begin{tabular}{lrrrr}
\toprule
shape & \texttt{alloced\_size} & \texttt{sparse\_bytes} & difference & \texttt{st\_blocks * 512} \\
\midrule
1 MiB HEAD/TAIL hole          & 1{,}056{,}768  & 1{,}032{,}192  &    24{,}576 &    24{,}576 \\
10 MiB hole                   & 10{,}485{,}760 & 10{,}452{,}992 &    32{,}768 &    32{,}768 \\
50 MiB hole                   & 52{,}428{,}800 & 52{,}396{,}032 &    32{,}768 &    32{,}768 \\
4 KiB-data-per-64 KiB chunked & 2{,}097{,}152  & 1{,}572{,}864  &   524{,}288 &   524{,}288 \\
\bottomrule
\end{tabular}
\end{center}

The branch now emits
\texttt{Some(alloced\_size.saturating\_sub(sparse\_bytes))}, with
saturation guarding the (not observed in practice but algebraically
possible) case where the kernel's residue exceeds the dstream-
reported allocation.

\subsection{The decmpfs Lift (EX-26 Hypothesis F)}
\label{ssec:decmpfs-lift}

EX-22 left decmpfs fail-closed because none of the open-source
readers agreed on the macOS-oracle number and the EX-22 fixture
carried only one decmpfs case (the small-payload xattr-stream-
stored shape). EX-26 added a second decmpfs variant (a 256~KiB
compressible payload that ditto routes into the resource-fork-
stored shape) and observed that both shapes are captured by a
single rule: the allocated bytes are the sum of any stream-backed
xattr \texttt{alloced\_size} values plus the (typically absent)
primary dstream \texttt{alloced\_size}.

\textbf{Observation.} On the EX-26 fixture the two decmpfs files
carry:

\begin{itemize}
  \item \texttt{compressed.txt} (52~KiB compressible JSON): no
    primary dstream; \texttt{com.apple.decmpfs} is
    \emph{stream-backed} (flags = \texttt{0x01}) with
    \texttt{stream\_dstream.alloced\_size = 4096};
    \texttt{com.apple.ResourceFork} is stream-backed with
    \texttt{stream\_dstream.alloced\_size = 0} (empty fork).
    Sum: 4096. Oracle \texttt{st\_blocks * 512 = 4096}.
  \item \texttt{compressed-big.bin} (256~KiB compressible binary):
    no primary dstream; \texttt{com.apple.decmpfs} is
    \emph{embedded} (flags = \texttt{0x02}) with payload
    \texttt{fpmc} magic, \texttt{compression\_type = 8} (lzvn
    resource-fork-stored), \texttt{uncompressed\_size = 262144};
    \texttt{com.apple.ResourceFork} is stream-backed with
    \texttt{stream\_dstream.alloced\_size = 4096}. Sum: 4096.
    Oracle \texttt{st\_blocks * 512 = 4096}.
\end{itemize}

The on-disk reality is that the \texttt{compression\_type} byte in
the decmpfs header tells you \emph{where} the compressed bytes
live (xattr inline vs.\ resource fork), but the bytes themselves
are always counted by summing the \texttt{stream\_dstream.alloced\_size}
fields of every xattr that is stream-backed. Embedded xattrs
contribute zero because their bytes live inline in the xattr
payload, which is part of the inode body's storage and does not
allocate separate extents.

The disagreement between \texttt{linux-apfs-rw} (publishes
uncompressed size / 512) and \texttt{apfs-fuse} (publishes zero)
that EX-22 documented is now resolved against the macOS oracle:
neither was right; the sum-of-stream-dstreams formula is the
public on-disk rule that reproduces \texttt{st\_blocks * 512}
exactly.

A residual case the EX-26 fixture did not exercise is the
\emph{incompressible-payload} shape, where ditto leaves
\texttt{UF\_COMPRESSED} off and emits no \texttt{com.apple.decmpfs}
xattr at all. Such files fall through to branch~1 (regular file
with a primary dstream) and emit \texttt{Some(alloced\_size)}
under the EX-22 baseline; on the EX-26 fixture the
\texttt{compressed-random.bin} case did exactly this.

\subsection{Per-Inode Results}

\textbf{Observation.} On the EX-26 fixture, all ten rows match the
\texttt{st\_blocks * 512} oracle exactly:

\begin{center}
\begin{tabular}{lllrr}
\toprule
path & branch & rule & picked & st\_blocks * 512 \\
\midrule
ordinary.txt          & 1 & \texttt{alloced\_size}                                              &      4096 &      4096 \\
clone.txt             & 1 & \texttt{alloced\_size}                                              &      4096 &      4096 \\
compressed-random.bin & 1 & \texttt{alloced\_size}                                              &  131{,}072 &  131{,}072 \\
sparse.bin            & 2 & \texttt{alloced\_size - sparse\_bytes}                              &     24{,}576 &     24{,}576 \\
sparse-medium.bin     & 2 & \texttt{alloced\_size - sparse\_bytes}                              &     32{,}768 &     32{,}768 \\
sparse-large.bin      & 2 & \texttt{alloced\_size - sparse\_bytes}                              &     32{,}768 &     32{,}768 \\
sparse-chunked.bin    & 2 & \texttt{alloced\_size - sparse\_bytes}                              &    524{,}288 &    524{,}288 \\
compressed.txt        & 3 & \texttt{decmpfs.stream\_dstream.alloced} (xattr-stream-stored)      &      4096 &      4096 \\
compressed-big.bin    & 3 & \texttt{ResourceFork.stream\_dstream.alloced} (resource-fork-stored)&      4096 &      4096 \\
link.txt              & 4 & \texttt{zero} (symlink)                                              &         0 &         0 \\
\bottomrule
\end{tabular}
\end{center}

Zero mismatches; every row picks a value the rule emits, against
the macOS oracle. The hard-linked file \texttt{hard.txt} does not
appear as a separate row because it shares its inode with
\texttt{ordinary.txt} and contributes through the same per-inode
rule.

\section{Directory Aggregates}

\textbf{A directory's total is the \emph{unique-inode logical sum
within the aggregate root}.}

A regular file's logical size contributes to a directory total once,
regardless of how many hard-linked paths within that directory tree
reference the same inode. A symlink contributes its target byte
length to its own row but contributes nothing to any directory total
(the target may live elsewhere or not exist at all, and even when it
exists its size is already counted at the target's row).

\textbf{Why this policy.} The naive policy --- sum the logical size of
every path under a directory --- catastrophically over-counts on real
directories with hard-linked content: mail spools, package caches,
Git object stores, snapshot archives. The user looking at a 12 GB
directory deserves to know that 11 GB of it is a single 11 GB file
that also appears elsewhere, not that the directory legitimately
holds 12 GB of unique content. WizTree-on-NTFS rarely encounters
hard-link density that exposes this; on APFS the case is common.

\textbf{Consequence.} Sibling directory totals are not strictly
additive into their parent when the same file identity appears in
more than one subtree. This is correct, not buggy, but the user
interface must label aggregate semantics so the user understands
what the number means.

The aggregate rule is implemented in one place
(\texttt{aggregate.py}, ported to the Rust crate's
\texttt{namespace} module) using
\texttt{setdefault(file\_id, logical\_size)} keyed by inode --- the
policy is enforced by construction rather than by careful summing.

Allocated-size aggregates use the same unique-inode rule, with the
extra step that any contributing inode whose
\texttt{allocated\_size} is \texttt{None} (sparse or decmpfs)
propagates \texttt{None} up to the directory total.

\section{What This Chapter Does Not Promise}

Three metrics remain outside the scope this chapter defends. Each
has its own oracle and its own corpus of edge cases; chapter~13
records them as open questions for future work.

\begin{description}
  \item[Exclusive size.] Bytes owned only by this inode, after
    accounting for clones and snapshots. Requires the
    extent-reference tree to be opened and joined against per-inode
    extents.
  \item[Shared size.] Bytes shared with at least one other inode or
    snapshot.
  \item[Snapshot-retained bytes.] Bytes that exist on disk solely
    because a snapshot still references them.
\end{description}

A future version of this manual will have its own chapter on each.
Until those chapters exist, the parser does not emit those columns.
Logical size is what it knows; allocated size is what it knows for
five-of-six case classes; sparse and decmpfs allocated sizes are
what it refuses to claim.
